% Document uses 12 pt font
% 1 in margins
% Contains a relative path for images

\documentclass [11pt]{article}

% page geometry 
\usepackage[margin=1in]{geometry}


% ----------  PACKAGES START ------------ %

% Table cell color package and highlighting
\usepackage[table]{xcolor}
\usepackage{color,soul}


% VIC title package
\usepackage{cabin}
\usepackage[T1]{fontenc}

% default font package
%\usepackage{times}
\usepackage{helvet}
\renewcommand{\familydefault}{\sfdefault}

% ---------- End Font Packages -------------- %

% Title Packages
\usepackage{titlesec}
\usepackage{titletoc}

% Image Package
\usepackage{graphicx}

% Table Packages
\usepackage{longtable}
\usepackage{multirow}
\usepackage{multicol}
\usepackage{multirow}
\usepackage{array}
\renewcommand{\arraystretch}{1.4}% Spread rows out evenly in table

% Color Packages
\usepackage{color}   
\definecolor{sectionC}{rgb}{0.016,0.227,.365}
\definecolor{subsectionC}{rgb}{.87,0.87,.87}
\definecolor{subsubsectionC}{rgb}{.94,.93,.90}
\definecolor{tableCell}{rgb}{.96,.95,.90}


% List package
\usepackage{enumitem}
\setenumerate{itemsep=0pt, itemindent=0in,leftmargin=0.5in}

% Paragraph parameter

\setlength{\parindent}{0pt}


% ------------- Creates a linked Table of Contents  Start --------------- %
\usepackage{hyperref}
\hypersetup{
colorlinks=true, %set true if you want colored links
linktoc=all,     %set to all if you want both sections and subsections linked
linkcolor=black,}  %choose some color if you want links to stand out

% ------------- Creates a click-able Table of Contents  End--------------- %

% ---------- PACKAGES END ------------ %



% ------------------- START HEADER AND FOOTER ---------------------------%
\usepackage{fancyhdr}

% Helps with the n of total n pages
\usepackage{lastpage}

\pagestyle{fancy}

% Header
\lhead{Proof of Concept }
\rhead{Revision: 0}
\fancyhead[LE,CO]{VIC - Group 6}

% Removes line under the header 
\renewcommand{\headrulewidth}{0pt}
\setlength{\headsep}{.2in}

% Footer 

% Set the right side of the footer to be the page number
\fancyfoot[R]{Page \textbf{\thepage}\ of \textbf{\pageref{LastPage}}}
\fancyfoot[C]{}

% ------------------- END HEADER AND FOOTER ---------------------------%




% -------- SECTION AND SUBSECTION FORMATING START -------- % 
% starts the 
%\setcounter{section}{1}


\titleformat{\section} % Section
{\normalfont \fontsize{14}{14} \bfseries}{}{0em}{\colorsection}

% Makes a background color
\newcommand{\colorsection}[1]{%
  \colorbox{sectionC}{\parbox{\dimexpr\textwidth-1\fboxsep}{\color{white}\Large\thesection\ \hspace{1mm} #1}}}

% Makes a background color
\titleformat{\subsection} % Subsection
{\normalfont \fontsize{12}{12}  \bfseries}{}{0em}{\colorsubsection }

\newcommand{\colorsubsection}[1]{%
  \colorbox{subsectionC}{\parbox{\dimexpr \textwidth -1\fboxsep}{\large\thesubsection\ #1}}}


% Makes a background color
\titleformat{\subsubsection} % Subsubsection
{\normalfont \fontsize{12}{12} \bfseries}{}{0em}{\colorsubsubsection}

\newcommand{\colorsubsubsection}[1]{%
  \colorbox{subsubsectionC}{\parbox{\dimexpr\textwidth-1\fboxsep}{\thesubsubsection\ #1}}}

% -------- SECTION AND SUBSECTION FORMATING END -------- % 
\usepackage{lipsum}


% -----  IMAGE PATH START -----%
% Relative Image Path
\graphicspath {figures/}
% -----  IMAGE PATH END -----%

% ------ PARAGRAPH FORMAT START ----%
%\setlength{\parskip}{.2em}% Sets the space between new paragraph items 
\setlength{\parindent}{0em} % paragraph indent
% ------ PARAGRAPH FORMAT END ----%




%------------------------------TOC FORMAT START --------------------------------%
\usepackage{tocloft}

% Section indentations
\cftsetindents{section}{0em}{1.5em}
\cftsetindents{subsection}{1em}{2em}
\cftsetindents{subsubsection}{2em}{3em}

% Toc title size
\renewcommand\cfttoctitlefont{\Large\bfseries}

% Removes bold headings from toc
%\renewcommand{\cftsecfont}{\normalfont}

% Removes bold heading page numbers from toc
\renewcommand{\cftsecpagefont}{\normalfont}

% add dots after headings
%\renewcommand{\cftsecleader}{\cftdotfill{\cftdotsep}} 


% number of section headings we want to see in toc
\setcounter{tocdepth}{2}

% Spaceing before headings in toc
\setlength{\cftbeforesecskip}{6pt}

% ------------------------------TOC FORMAT END --------------------------------%








% -------------- DOCUMENT START ---------------%
\begin{document}

% --------- TITLE PAGE START ------- %
\begin {center} 

\thispagestyle{empty}
\vspace*{5cm}

% Logo Insertion
\begin {figure}[h!]
\centering
\hspace{-10mm}\includegraphics [scale = .5, trim={.4cm 0 .8cm 0},clip] {figures/vic_logo.png}
\end {figure}

{\fontfamily{\cabinfamily}\selectfont
\Huge{Vehicle Intersection Control} }

\vspace{1 cm}
{\Large{\textsc{McMaster University}}\\}  \vspace {1cm}
{\large Proof of Concept Demonstration\\ \vspace {0.4 cm} SE 4G06}  \vspace {1cm}

{\large \textsc{Group 6} \\} \vspace{1cm}

{
Alex Jackson \\
Jean Lucas Ferreira \\
Justin Kapinski\\
Mathew Hobers\\
Radhika Sharma\\
Zachary Bazen}




\end{center}

% --------- TITLE PAGE END------- %

\pagebreak

% Inserting table of contents and table of figures 

\tableofcontents
%\listoftables
% \listoffigures



\pagebreak

% -----------  REVISION HISTORY START ----------- %

%%\section*{Revisions}
%\thispagestyle{empty}
%\section{Revisions}
%\begin{longtable}{| p{.2\textwidth } | p{.23\textwidth } | p{.23\textwidth } | p{.23\textwidth } |}
%
%\hline 
%\centering \textbf{Date} & 
%\multicolumn{1}{c}{\textbf {Revision Number}} &
%\multicolumn{1}{|c}{\textbf {Authors}} & 
%\multicolumn{1}{|c|}{\textbf {Comments}} \\ \hline
%
%\multirow{4}{*}{\centering December 1, 2016}  & 
%\multirow{4}{*}{Revision 0}& 
%{Alex Jackson \newline
%Jean Lucas Ferreira \newline
%Justin Kapinski\newline
%Mathew Hobers\newline
%Radhika Sharma\newline
%Zachary Bazen}
%&
% \multirow{4}{*}{N/A} \\ 
%\hline 
%
%\caption{VIC Table of Revisions} 
%\end{longtable}
%% -----------  REVISION HISTORY END ----------- %
\pagebreak


\section{Challenges}

\subsection{Lane Following}

 Follow existing software solutions and modify it to suit the needs of the project requirements \\ \\
To capture images in real time we plan to attach a webcam to each vehicle.  \\
Proceeding through the intersection: we intend the vehicle to proceed at slow speeds in a straight line until sufficient data is available to follow lanes again

%- Use existing software or make own \\
%- look at previous examples to build own algorithm

\subsection{Car to Intersection Controller Communication}

 Vehicles will communicate to the intersection controller via Bluetooth.

\subsection{Overall Power Supply}
 Each vehicle will be fitted with various sensors, microcontrollers, and a webcam. All of these electronic components will be a strain on the power supply. Thus,  we will buy an additional power supply to be used while the other is charging. 

\subsection{How far the car as traveled}
Knowing the distance each car has travelled will allow us to perform various calculations, such as velocity control. \\
This will be accomplished by implementing a Hall effect sensor on each car.


\subsection{Obstacle Detection}
Obstacles may be present on track, and collision with such obstacles must be avoided. \\
We can define an obstacle as any foreign object that cannot be safely driven over by the car. \\
Obstacle detection will be possible through a sensor in front of the car. The sensor will be limited to detect obstacles only within the current driving lane. A viable sensor to implement is an infrared sensor. 

% - Possible obstacles include cardboard cutouts, wooden blocks, and other cars using the track


\section {Software Challenges}
\subsection{Processing speed}
Some code will be critical to the performance and it might not be possible to write this code in a high level language like Python. \\We will consult how past groups dealt with this problem, and consider an optimal language to satisfy the processing requirements of the vehicle.
 

\section{Open Issues}


\subsection{Intersection Computer and Microcontroller}

\begin{itemize}
    \item  What programming languages should be used and how to interface between them
    \item How will the algorithms access the data from the sensors
    \item Will the processing power of the Raspberry Pi be good enough to run image processing algorithms in real time
\end{itemize}



\end{document}
